\section{来源审计：从“被动小猫”到开放世界智能体}

这堂课的核心问题不是“如何再做一个更强的 Minecraft bot”，而是如何把 foundation model 的广泛先验接入一个能够\textbf{观察、行动、接受后果并持续积累技能}的闭环。Jim Fan 先用经典视觉发展实验说明主动经验的重要性，再把 MineDojo、MineCLIP、Voyager、Eureka 与 VIMA 串成一条系统路线：开放环境提供无穷任务，互联网提供人类知识，foundation model 提供可迁移表示，而 reward、code、memory、curriculum 与 simulator feedback 让这些先验真正进入行动。

\subsection{官方课堂边界与证据层级}

本次重写以 Stanford Online 官方录像 \texttt{wwQ1LQA3RCU} 为课堂顺序来源，课堂日期为 2023 年 10 月 24 日，录像时长 47 分 48 秒，并提供人工 \texttt{en-US} 字幕。61 个教学状态从 1080p 录像中恢复，覆盖所有实际承担概念、架构、结果或例子作用的页面；重复播放帧与纯过渡帧不重复计入。论文只用于核对术语与机制，不能反过来把 2023 年 10 月之后的结果写成课堂已知事实。

\begin{longtable}{@{}L{0.18\textwidth}L{0.31\textwidth}L{0.40\textwidth}@{}}
\toprule
证据层 & 本地材料 & 在讲义中的职责 \\
\midrule
官方录像 & 47:48、61 个教学状态 & 决定讲授顺序、视觉证据、demo 与课堂结论边界 \\
官方人工字幕 & 985 条 caption、timed transcript & 恢复动机、警告、比喻、限制条件和章节转场 \\
Primary sources & MineDojo、Voyager、Eureka、VIMA、VPT、RT-2、RoboCat & 核对训练目标、架构名称、实验对象与公式含义 \\
覆盖与视觉审计 & manifest、selection TSV、coverage matrix、PDF QA & 证明每张教学图被使用且成品可读，不以“文件存在”替代验收 \\
\bottomrule
\end{longtable}

\begin{warningbox}{时间边界与证据边界}
本文讲的是 2023 年课堂里的 generalist-agent research agenda。MineDojo、Voyager 或 Eureka 在某个 simulator 和指标上取得优势，不等于已经得到跨环境、跨 embodiment 的通用智能；VIMA、RT-2、RoboCat 展示多模态与机器人迁移，也不等于解决了安全、长期记忆、开放世界评价或真实世界数据瓶颈。
\end{warningbox}

\subsection{Held--Hein 实验：相同视觉输入为何产生不同能力}

课程从 Held 与 Hein 的“两只小猫”实验切入，是为了把 embodiment 从“拥有一副机器人身体”提升为更严格的学习条件：\term{embodiment（具身）}指智能体的感知与动作通过身体和环境动力学形成闭环，动作会改变下一时刻看到什么，也会产生可归因的成功、失败与代价。主动小猫控制转盘运动，被动小猫获得近似匹配的视觉流，却没有相同的行动因果关系。

\lecturefigure{slide-01-two-kittens.jpg}{Held--Hein 两只小猫实验：主动行动与被动观看}{官方录像恢复幻灯片 V001；00:03:18}

\begin{knowledgebox}{读图：不要只比较“看了多少”}
图中两只小猫获得相近的视觉刺激，但左侧被动个体的运动由右侧主动个体驱动。先比较\textbf{谁决定动作}，再比较视觉流；关键差异不是像素数量，而是 action 与 observation 之间是否存在可学习的因果对应。实验支持主动 sensorimotor experience 很重要，但它本身不能证明所有认知能力都必须从零开始通过 physical interaction 获得。
\end{knowledgebox}

智能体与环境的最小闭环可以写成：
\[
o_t = \mathcal{O}(s_t),\qquad a_t \sim \pi_\theta(\cdot\mid h_t),\qquad
s_{t+1}\sim P(\cdot\mid s_t,a_t),\qquad h_{t+1}=f(h_t,o_{t+1},r_t).
\]
其中 $s_t$ 是环境真实状态，$o_t$ 是可见 observation，$a_t$ 是 action，$\pi_\theta$ 是参数为 $\theta$ 的 policy，$P$ 是环境动力学，$r_t$ 是反馈或 reward，$h_t$ 是由历史构成的内部状态。被动观看只提供部分 $o_t$；主动学习还拥有“我执行了哪个 $a_t$，为何得到 $o_{t+1}$ 与 $r_t$”这一因果链。

\teachervoice{讲者把 ChatGPT 称为极其强大的“passive kitten”：它从互联网吸收了巨量人类知识，却主要以观察者身份学习。课堂主张不是丢掉 pretraining、让机器人从零探索，而是在强大的被动先验之上增加主动闭环，使模型学会行动后果、可供性与长期任务。}

\begin{importantbox}{本讲的起点不是 pretraining 与 embodiment 二选一}
互联网预训练解决“人类已经知道什么”，主动交互解决“我的动作在当前身体与环境中会造成什么”。Generalist agent 需要把二者组合：用 foundation model 提供 broad prior，用环境反馈校准 grounding、skill 与 control。
\end{importantbox}

\subsection{本章小结}

本章确定了来源、时间与证据边界，并用两只小猫实验建立全讲主张：相同 observation 并不等于相同 learning signal，行动控制权与反馈闭环不可缺失。后文所有系统都可以看作在回答同一个问题：如何把 passive internet knowledge 转换为可执行、可验证、可积累的 active competence。

